\section{多 UV 面的算术修复与可复核执行路径}
\label{sec:precision-repair}

从图构造出正确的递归表达式之后，还必须可靠地计算其完整系数和。
本节记录四圈实际反例带来的修复。Bogoliubov BPH 递归、完整零动量
Taylor 条件及共同原始边参数均保持不变；新增机制处理的是有限表达式的算术，
不另选物理方案，也不调用逐拓扑预写的生产积分公式。

\subsection{为什么单面阈值遗漏了实际四圈错误}

\texttt{L4E3G205} 有三个顶点不交的二点 bubble 子图，原始边集分别为
$\{0,1\}$、$\{3,4\}$、$\{6,7\}$。令每个子图中的两条边各为
$\delta/2$，其余六条边各为 $(1-3\delta)/6$，并取
$M=p_1^2=p_2^2=p_3^2=1$。下列点都位于参数单纯形内部，
各子图参数和及补集参数和均大于原来的 $10^{-4}$ 阈值：
\begin{center}\small
\begin{tabular}{rll}\toprule
$\delta$ & 绕过保护的 binary64 密度 & 80/110 位参考及修复值\\\midrule
$10^{-2}$ & $-99552$ & $-99369.9951744098$\\
$10^{-3}$ & $-412316860416$ & $-81064258.2517440$\\\bottomrule
\end{tabular}
\end{center}
第一组相对误差约为 $0.183\%$，第二组相对误差约为 $5085$ 倍。
多个二阶 Taylor 减除同时作用时，共同抵消可以远强于每个单独面尺度所暗示的程度。
加大采样数不能消除这种系统算术误差。
不受保护的计算只用于重现历史缺陷，不作为被接受的积分密度。

\subsection{完整 Taylor 系数和的相对精度触发}

记第 $W$ 个递归项中保留的 Taylor 系数为
\[
 a_{W,\mathbf k}(\alpha)
 =[t^{\mathbf k}]\,U_W(t)^{-3}\Phi_E(H_W(t),M).
\]
除完整密度外，程序同时计算所有保留系数的绝对值之和：
\[
 \mathcal K_G=\sum_W(-1)^{|W|}\sum_{\mathbf k}a_{W,\mathbf k},
 \qquad
 A_G=\sum_W\sum_{\mathbf k}|a_{W,\mathbf k}|.
\]
原有 UV 面检查继续保留。对未被单面检查选中的点，若
\begin{equation}
 64\,\varepsilon_{\rm work}\,A_G
 >10^{-10}|\mathcal K_G|,
 \label{eq:relative-cancellation-trigger}
\end{equation}
或浮点计算得到非有限结果，就以同一输入重新计算完整递归和。
这里的绝对触发目标为零；$\varepsilon_{\rm work}$ 是实际工作类型的机器精度。
批量路径记录 \texttt{numpy.longdouble} 的真实尾数位数，
本机仍只有 53 个有效二进制位，不能从类型名称推定额外精度。

取同一 $\delta=10^{-3}$ 参数点，但令 $p_i^2=10^{-28}$，
未保护的值约为 $-4.1633\times10^{-17}$，而 110 位参考为
$-8.7361\times10^{-21}$。
若使用 $10^{-12}$ 的固定绝对阈值，该错误会因数值很小而被放过。
式~\eqref{eq:relative-cancellation-trigger} 会触发这个点，并恢复其相对精度。

$A_G$ 本身由有限精度算术求得，Jet 卷积和 Schur 消元的内部误差
没有被转换为严格的向外舍入区间。因此该规则是保守的数值触发条件，
并不是已经证明的全域向前误差界。
实际回归、未触发点的复查与独立解析参考继续分别记录。

\subsection{标准库复算与动态工作精度}

标准库 Decimal 保留完整且独立的执行路径。
把 $1$ 也计入各真子图及其补集参数和，记其中最小正数为 $r$，初始位数取
\[
 d=\min\{470,\max\{80,
           32+12\lceil\max(0,-\log_{10}r)\rceil\}\}.
\]
在该位数计算 $v_d$ 后，再以 $d'=\min(500,d+30)$ 位计算，要求
\[
 |v_{d'}-v_d|
 \le10^{-25}\max\{|v_{d'}|,|v_d|\}.
\]
两次均为精确零时接受；其余情况使用相对复核。
若不稳定则继续提高位数，最高 500 位仍未稳定时明确报错。
小参数的 $\log(1+z)$ 与 $\Phi_2$ 继续按当前上下文动态求和，
保留上一节两个真实历史反例的回归，避免固定项数级数在不同位数下共享同一偏差。
高精度算术使用已接收到的 binary64 输入及 Gram 数据，
不会增加原始输入本身的有效位数。

\subsection{可选原生执行器以 50/100 位逐点复核}

Python 从图生成子图、递归项和相容路由；可选 C++17 执行器只读取
生成程序中的整数路由系数、缩放单项式、Taylor 阶数及递归符号。
它不按图号选择拓扑公式。
Boost.Multiprecision 分别以 50 位和 100 位重新执行同一程序。
对非零结果，其接受检查包括
\[
 \max\left\{
  \frac{|v_{100}-v_{50}|}{|v_{100}|},\,
  \frac{10^{-90}A_{100}}{|v_{100}|}
 \right\}\le10^{-25}.
\]
零值另作明确检查。原生精度检查失败时，交给自适应 Decimal，
不静默采用未经确认的低精度值。
这同样是数值检查，不是严格舍入误差证明。

库加载器核对 ABI、操作系统与 CPU、当前 C++ 源码的 SHA256 和库文件的 SHA256；
缓存句柄在后续调用也重新核查。
库缺失、过期或属于不同主机时，标准库路径仍可工作。
\texttt{PHI3LAB\_DISABLE\_NATIVE=1} 可明确禁用原生后端。
该库由本项目随包源码构建，与独立参考所用的官方 pySecDec 库分别记录。

多精度结果输出到 binary64 时可能下溢，状态单独报告；
100 位文本接口保留未转为 binary64 的数值。
不能对 binary64 无法表示的非零结果承诺该输出类型的相对精度。

\subsection{实际验证和执行成本}

本次冻结源码下实际完成 15 个共同面、不等面、边和顶点重标号、
真嵌套及极小不变量回归，并另外检查 20 个实际未触发点。
80/110 位 Decimal、标量、批量和原生 50/100 位检查全部通过；
最大标量相对差约为 $2.17\times10^{-13}$，批量约为 $1.83\times10^{-13}$。
完整输入、程序哈希、库身份及触发计数保存在
\path{results/automatic_multiface_comparison.json}。

完整标准库套件在禁用 site packages 且设置
\texttt{PHI3LAB\_DISABLE\_NATIVE=1} 后实际执行 170 项，
明确跳过 34 项可选测试，全部执行项通过。
可选 NumPy、SciPy 和原生环境另实际执行 85 项，无跳过、失败或错误。
两个套件含重叠测试，计数不相加。
另外明确禁用原生后端的代表检查实际通过 8 项，跳过 3 项 NumPy 批量测试。
各份记录包含源码指纹及日志哈希。

1024 点实测基准包含 32 个共同 UV 面点和 992 个 Dirichlet $\beta=0.35$ 点。
110 位 Python Decimal 实际耗时 18.9802 秒；
同一生成程序的原生 1、2、4 线程分别为 3.5847、1.8712、0.99724 秒。
4 线程约快 19.03 倍，最大相对差为 $2.09\times10^{-16}$。
这些数字描述指定机器、输入和工作位数下的算术成本，
不能直接换算成所有完整积分的加速比。
原始每点记录和时间保存在 \path{docs/automatic-native-benchmark.json}；
冻结后的数学依赖哈希核对没有被称为重新计时。

本节的点态修复、完整振幅求积和解析收敛证书仍分别记录。
新增高精度执行路径使暴露的问题被实际重算和回归检查，
但点态数值触发、随机组误差及确定性求积细化都不因此成为严格的总积分误差界。
